%! AP Statistics — Unit 4, Lesson 4.1 — Guided Notes KEY
\documentclass[10pt]{article}
\usepackage{apstats-article}
\usepackage{apstats-key}

\begin{document}

\pageheader{Unit 4, Lesson 4.1}{Guided Notes --- KEY}
\namedateperiod

\vspace{0.12in}

\begin{objectivebox}
\small By the end of this lesson, I will be able to\dots\\[4pt]
\ans{Identify patterns in data and write statistical questions those patterns suggest.}\\[1pt]
\ans{Explain why a pattern alone does not confirm a non-random process.}\\[1pt]
\ans{Recognize that random processes can produce streaks, clusters, and other surprising patterns.}
\end{objectivebox}

\vspace{0.1in}

\begin{vocabbox}
\termblanklong{Random process} \ans{A process whose outcome cannot be predicted with certainty; determined by chance.}
\termblanklong{Pattern} \ans{An apparent regularity or structure observed in data.}
\termblanklong{Streak} \ans{A run of identical consecutive outcomes in a sequence.}
\termblanklong{Random variation} \ans{Natural fluctuation in data produced by chance alone.}
\termblanklong{Statistical question} \ans{A question that anticipates variability in the data needed to answer it.}
\end{vocabbox}

\vspace{0.1in}

\begin{hookbox}
\small\textbf{Coin flip sequences:}
\begin{itemize}[leftmargin=1.5em, itemsep=3pt]
  \item Sequence A: H H T H T T H H H T T H H T H T A T H T H
  \item Sequence B: H H H H H T T T T T H H H H H T T T T T
  \item Sequence C: H T H T T H T H H H T T H T H T H H T H
\end{itemize}
\textbf{Which looks most like a real fair coin? Which looks suspicious? Why?}\\[4pt]
\ans{Sequences A and C look most like real coin flips because they have both runs and alternations mixed together. Sequence B looks suspicious because the long blocks of 5 identical outcomes seem too regular --- but in fact, even Sequence B can arise from a fair coin; our intuition tends to underestimate how often runs appear.}
\end{hookbox}

\vspace{0.1in}

\begin{notesbox}{Patterns in Random Processes}
\small
\textbf{Key Principle (VAR-1.F.1):} Patterns in data \ans{do not} necessarily mean
that the variation is \ans{non-random}.

\vspace{10pt}
\renewcommand{\arraystretch}{1.6}
\begin{tabularx}{\linewidth}{@{} >{\bfseries\color{navy}}p{0.3\linewidth} X @{}}
Random process & \ans{A process whose outcome is determined by chance; outcomes cannot be perfectly predicted.} \\
Pattern & \ans{An apparent regularity in data --- a streak, cluster, or trend that catches attention.} \\
Streak & \ans{A run of identical consecutive outcomes (e.g., 6 heads in a row).} \\
\end{tabularx}

\vspace{10pt}
\textbf{The Hot-Hand Example:}\\
A basketball player makes 7 consecutive free throws (career average: 70\%).
\begin{itemize}[itemsep=4pt]
  \item The pattern suggests: \ans{She is ``on fire'' / has a hot hand right now.}
  \item But we should ask: \ans{How likely is a streak of 7 for a 70\% shooter even by chance?}
  \item To answer that question, we need: \ans{Probability.}
\end{itemize}

\vspace{8pt}
\textbf{Counter-intuitive fact:} Most people expect random sequences to have
\ans{more} alternations. In reality, a run of 5+ identical outcomes in
20 coin flips has about a \ans{12\%} chance of occurring.
\end{notesbox}

\vspace{0.1in}

\begin{notesbox}{Identifying Statistical Questions from Patterns}
\small
A \textbf{statistical question} is one that \ans{anticipates variability in the data needed to answer it}.

\vspace{8pt}
For any pattern in data, ask these questions:
\vspace{4pt}
\renewcommand{\arraystretch}{1.8}
\begin{tabularx}{\linewidth}{@{} l X @{}}
\textbf{Question} & \textbf{My answer} \\
\hline
What pattern do I observe? & \ans{Identify the streak, cluster, or trend specifically.} \\
What claim is the pattern tempting me to make? & \ans{Name the non-random cause being implied (hot hand, environmental hazard, etc.).} \\
Could chance alone produce this pattern? & \ans{Consider how surprising the pattern is given a random model.} \\
What evidence would convince me it is NOT random? & \ans{Many repeated observations; controlled study; probability calculation showing the pattern is very unlikely by chance.} \\
\end{tabularx}
\end{notesbox}

\vspace{0.1in}

\begin{practicebox}
\small

\vspace{6pt}
\begin{enumerate}[leftmargin=1.5em, itemsep=12pt]
  \item A small town reports 8 cases of a rare disease in one year when the expected rate
        predicts fewer than 2.
        \begin{enumerate}[label=(\alph*), itemsep=4pt]
          \item \ans{A cluster of 8 cases where fewer than 2 are expected.}
          \item \ans{Is the rate of this disease in this town higher than the national average, or could the cluster be due to random variation in a small population?}
        \end{enumerate}

  \item A coffee shop notices that customers order hot drinks on 9 of the last 10 rainy days.
        \begin{enumerate}[label=(\alph*), itemsep=4pt]
          \item \ans{Hot drink orders appear much more common on rainy days.}
          \item \ans{Is the proportion of hot drink orders higher on rainy days than on non-rainy days, or could this association arise by chance?}
        \end{enumerate}

  \item A student flips a coin 15 times and gets heads 11 times.
        \begin{enumerate}[label=(\alph*), itemsep=4pt]
          \item \ans{11 heads out of 15 flips (73\%) --- more heads than expected from a fair coin (50\%).}
          \item \ans{Yes --- a fair coin produces 11 or more heads in 15 flips about 5.9\% of the time, which is unusual but not impossible.}
        \end{enumerate}
\end{enumerate}
\end{practicebox}

\end{document}
